% ============================================
% DAFTAR SIMBOL
% File: Simbol.tex
% ============================================

% Gunakan package glossaries untuk simbol

% ============================================
% 1. SIMBOL GRAF & STRUKTUR DATA
% ============================================

\newglossaryentry{sym:G}{
    type=symbols,
    name={\ensuremath{G}},
    description={\textit{Knowledge Graph} yang didefinisikan sebagai $G = (E, R)$},
    sort=G
}

\newglossaryentry{sym:E}{
    type=symbols,
    name={\ensuremath{E}},
    description={Himpunan entitas dalam \textit{Knowledge Graph}},
    sort=E
}

\newglossaryentry{sym:R}{
    type=symbols,
    name={\ensuremath{R}},
    description={Himpunan relasi dalam \textit{Knowledge Graph}},
    sort=R
}

\newglossaryentry{sym:T}{
    type=symbols,
    name={\ensuremath{T}},
    description={Himpunan \textit{triplet} pengetahuan $(h,r,t)$},
    sort=T
}

\newglossaryentry{sym:Nraw}{
    type=symbols,
    name={\ensuremath{N_{\text{raw}}}},
    description={Jumlah total data mentah publikasi yang dikumpulkan},
    sort=Nraw
}


% ============================================
% 2. SIMBOL ENTITAS & TRIPLET
% ============================================

\newglossaryentry{sym:h}{
    type=symbols,
    name={\ensuremath{h}},
    description={Entitas asal (\textit{head}) dalam \textit{triplet}},
    sort=h
}

\newglossaryentry{sym:r}{
    type=symbols,
    name={\ensuremath{r}},
    description={Relasi atau predikat dalam \textit{triplet}},
    sort=r
}

\newglossaryentry{sym:t}{
    type=symbols,
    name={\ensuremath{t}},
    description={Entitas tujuan (\textit{tail}) dalam \textit{triplet}},
    sort=t
}


% ============================================
% 3. SIMBOL SYSTEM FLOW (RAG)
% ============================================

\newglossaryentry{sym:q}{
    type=symbols,
    name={\ensuremath{q}},
    description={Pertanyaan atau kueri dari pengguna},
    sort=q
}

\newglossaryentry{sym:a}{
    type=symbols,
    name={\ensuremath{a}},
    description={Jawaban atau respons yang dihasilkan sistem},
    sort=a
}

\newglossaryentry{sym:Kclues}{
    type=symbols,
    name={\ensuremath{\mathcal{K}^{clues}}},
    description={Kata kunci petunjuk awal (\textit{semantic clues})},
    sort=Kclues
}

\newglossaryentry{sym:Kl}{
    type=symbols,
    name={\ensuremath{\mathcal{K}_l}},
    description={Kata kunci spesifik (\textit{Low-level Keywords})},
    sort=Kl
}

\newglossaryentry{sym:Kh}{
    type=symbols,
    name={\ensuremath{\mathcal{K}_h}},
    description={Kata kunci tema besar (\textit{High-level Keywords})},
    sort=Kh
}

\newglossaryentry{sym:Llocal}{
    type=symbols,
    name={\ensuremath{\mathcal{L}_{\text{local}}}},
    description={Konteks informasi lokal dari subgraf},
    sort=Llocal
}

\newglossaryentry{sym:Lglobal}{
    type=symbols,
    name={\ensuremath{\mathcal{L}_{\text{global}}}},
    description={Konteks informasi global dari agregasi relasi},
    sort=Lglobal
}


% ============================================
% 4. METRIK EVALUASI
% ============================================

\newglossaryentry{sym:absV}{
    type=symbols,
    name={\ensuremath{|V|}},
    description={Jumlah klaim yang terverifikasi (\textit{verified claims})},
    sort=V
}

\newglossaryentry{sym:absS}{
    type=symbols,
    name={\ensuremath{|S|}},
    description={Total jumlah klaim dalam jawaban (\textit{total statements})},
    sort=S
}

\newglossaryentry{sym:sim}{
    type=symbols,
    name={\ensuremath{\text{sim}(q, \hat{q})}},
    description={Fungsi kemiripan (\textit{similarity}) antara dua teks},
    sort=sim
}
